\begin{table*}[htbp]
\centering
\scriptsize
\caption{The five predictions of the experimental design, the criterion each was held to, the measured values on the committed bake or on the experiments of Section~\ref{sec:experimental-design}, and whether each held; H1 is the acceptance and refutation criterion fixed for the method before this evaluation and H2 the second clause of the research brief that proposed it.}
\label{tab:hypotheses}
\begin{tabular}{lp{4.8cm}p{6.2cm}l}
\toprule
Id & Prediction & Measured & Outcome \\
\midrule
H1 & \relax At a size of at most 1\%, the stale-volatility model is flagged in at least 30\% of replications; refuted below 10\% & \relax union 99.9 (99.4--100.0)\% at held-out size 0.9 (0.5--1.7)\%; regulatory zone 0.0\% & held \\
H2 & The Wald test flags wrong-correlation and missing-basis at least as often as the KS distance alone, and the no-deficiency desk at most at the stated rate & \relax wrong correlation 100.0\% against 100.0\%; missing basis 0.9\% against 86.7\%; no deficiency 0.8\% & held in part \\
H3 & \relax The false-alarm rate of the union stays at most 1\% on desks that differ from the calibration null & \relax 4.0\% (persistence 0.995), 6.2\% (dispersion 0.50), 99.0\% (pricing noise 30\%) & failed \\
H4 & The flag rate of the rolling-mean statistic rises with the age of the volatility input & \relax rates 5.7, 57.6, 97.5, 99.9, 100.0, 100.0\% at 2, 5, 10, 20, 40 and 80 days & held \\
H5 & The complement flags a constant mis-scale that the regulatory zone does not & \relax Wald 100.0\% against zone 0.0\% at a 10\% mis-scale & held \\
\bottomrule
\end{tabular}
\end{table*}
